\section{Post-training II: optimización de preferencias (DPO/PPO)}

\subsection{De la señal supervisada a la señal de preferencias}
Después del SFT, el modelo ya cuenta con la capacidad básica para ejecutar tareas, pero el estilo de las respuestas, la preferencia de riesgo, la concisión y la utilidad todavía necesitan moldearse más mediante el aprendizaje de preferencias. El expositor presentó numerosas observaciones de ingeniería en torno a las diferencias prácticas entre DPO y PPO.

\begin{importantbox}{Por qué la optimización de preferencias es efectiva para reasoning}
Las tareas de reasoning suelen tener múltiples soluciones viables; la optimización de preferencias puede guiar al modelo hacia una distribución de salidas más interpretable, con pasos más completos y errores más corregibles, en lugar de limitarse a igualar una única etiqueta supervisada.
\end{importantbox}

\begin{figure}[H]
\centering
\includegraphics[width=0.9\textwidth]{frame-05-dpo-ppo.jpg}
\caption{Video aproximadamente en 00:48:30: comparación de DPO/PPO y discusión sobre la estabilidad del entrenamiento}
\end{figure}

\subsection{Las disyuntivas de ingeniería entre DPO y PPO}
\begin{table}[H]
\centering
\renewcommand{\arraystretch}{1.2}
\begin{tabular}{p{2.4cm}p{4.0cm}p{3.8cm}p{3.0cm}}
\toprule
\textbf{Método} & \textbf{Ventajas} & \textbf{Riesgos} & \textbf{Etapa de aplicación} \\
\midrule
DPO & La implementación es relativamente simple, el costo es menor y la cadena de entrenamiento es corta & Es sensible a la calidad de las muestras de preferencias y requiere manejar el sesgo de longitud & Iteración rápida a gran escala \\
PPO/RL en línea & Puede aprovechar la retroalimentación en línea para optimizar de manera continua, con mayor capacidad de exploración & Es inestable, el ajuste de hiperparámetros es complejo y el costo es alto & Etapas de impulso final para capacidades clave \\
\bottomrule
\end{tabular}
\caption{Diferencias prácticas entre DPO y PPO}
\end{table}

\begin{knowledgebox}{Un tema que se repite a lo largo del curso}
Cambiar de algoritmo casi nunca es el factor decisivo; \textbf{la calidad de la señal de recompensa} y \textbf{la cobertura de la distribución de entrenamiento} suelen ser más importantes. Muchas ``ganancias algorítmicas'' en realidad provienen de correcciones en los datos y en la función objetivo.
\end{knowledgebox}

\begin{warningbox}{Dos tipos de degradación implícita en el aprendizaje de preferencias}
\begin{itemize}
    \item \textbf{Complacencia excesiva}: el modelo aprende a «parecer una buena respuesta», pero carece de profundidad de razonamiento real;
    \item \textbf{Atajo de longitud}: el modelo alarga o acorta la salida para activar la preferencia de reward, en lugar de mejorar la corrección.
\end{itemize}
\end{warningbox}

\subsection{Resumen de la sección}
La optimización de preferencias es la capa de control del comportamiento dentro del post-training. DPO y PPO no son una disyuntiva binaria, sino que deben desplegarse por etapas combinando la calidad de los datos, el presupuesto y los objetivos de estabilidad.

